\section{问题分析}
\subsection*{问题一：}
%也可以考虑使用\paragraph{问题一：} 不占用章节编号层级
\begin{enumerate}[label=(\arabic*), leftmargin=2.5em, itemsep=0.3em]
\item \textbf{问题分析是论文的重要部分。} 主要分析每个问题的特点、难点及建模思路，为后续模型建立提供依据。

\item \textbf{一般按问题分别分析。} 建议按照问题一、问题二、问题三依次展开，每一问对应一份分析。

\item \textbf{重点分析思路。} 说明为什么采用该模型、哪些因素需要重点考虑、模型如何解决问题，而不是直接推导公式。

\item \textbf{建立问题之间的联系。} 若后续问题依赖前一问的结果，应说明各问题之间的数据、参数或模型关系。

\item \textbf{不写求解结果。} 本部分只分析建模思路，不提前给出计算结果、实验结果或最终结论。
\end{enumerate}

\textbf{问题分析的目的，是告诉评阅老师："为什么这样建模"，而不是"模型怎么算"。}

\textbf{问题分析主要写什么？}

对于每一问，建议围绕以下几个方面展开：
\begin{enumerate}[label=(\arabic*), leftmargin=2.5em, itemsep=0.3em]
\item \textbf{明确问题类型} 判断属于评价、预测、分类、优化、决策、机理分析等哪一类问题。

\item \textbf{分析关键因素} 找出影响结果的主要变量、约束条件和题目提供的数据，说明哪些因素需要重点考虑。

\item \textbf{确定建模思路} 结合问题特点，分析适合采用哪些数学模型或算法，并简要说明选择理由。

\item \textbf{说明求解目标} 明确模型最终需要输出什么结果，如最优方案、预测值、评价结果、分类结果等。
\end{enumerate}

\textbf{写作提示}

\begin{enumerate}[label=(\arabic*), leftmargin=2.5em, itemsep=0.3em]
\item 先分析问题，再选择模型，不要为了使用某个模型而分析问题。

\item 模型名称不宜罗列过多，应选择最符合题意的方法，并说明选择依据。

\item 每一问尽量围绕"问题特点—关键因素—建模思路—预期结果"展开，保持逻辑清晰。

\item 问题分析可以在完成模型建立后再回过头完善，但内容必须与正文保持一致。
\end{enumerate}

\textbf{写作原则：分析问题、明确思路、说明方法、不写结果。}

\subsection*{问题二：}

\subsection*{问题三：}
